% =============================================================================
% Section 2 — Related Work
% =============================================================================

\section{Related Work}
\label{sec:related}

\subsection{AAC Systems for Neurodegenerative Disease}

Augmentative and alternative communication (AAC) for people with MND/ALS has been
studied extensively since the 1990s~\cite{beukelman2013aac}. Commercial solutions
dominate clinical practice: high-end eye-tracking devices such as Tobii Dynavox PCEye
and Eyegaze Edge offer gaze-to-grid selection with accuracy exceeding 98\,\% under
laboratory conditions, but carry hardware costs upward of USD~3{,}000 and require
periodic re-calibration by trained practitioners~\cite{beukelman2013aac}. These systems
presuppose sufficient residual oculomotor control, which deteriorates as MND progresses
toward ophthalmoplegia. Software AAC packages such as Grid~3 depend on stable Bluetooth
or Wi-Fi connectivity to synchronise symbol sets and user profiles, creating a~barrier
for home users without reliable broadband.

A survey by Clarke et al.~\cite{clarke2021aac} identified that AAC provision for
people with ALS is characterised by late referral, rapid device transitions as motor
function declines, and a~persistent gap in low-cost, adaptable solutions suitable for
the home environment. Brain-computer interfaces (BCI) via electroencephalography offer
an alternative pathway that bypasses the motor system entirely, but current non-invasive
BCI headsets achieve typical information transfer rates of 5--25 bits per minute in
optimal conditions, require trained operators, and are poorly suited to the home care
setting~\cite{zander2011passive}. Implanted cortical BCIs achieve higher rates but
carry surgical risk. NeuroAssist AI targets the intermediate stage — after limb and
speech function is lost but while a~degree of voluntary facial movement persists — for
which no low-cost, network-independent solution currently exists in the literature.

\subsection{Facial Gesture Detection via Computer Vision}

Three generations of camera-based facial analysis are relevant to this work.

\emph{First generation — cascade classifiers.} Viola and Jones~\cite{violajones2004}
introduced real-time face detection using Haar feature cascades, enabling deployment on
consumer hardware. These detectors localise the face but do not provide sub-facial
landmark coordinates; detecting blink requires additional heuristics applied to the
detected eye region and is highly sensitive to illumination variation.

\emph{Second generation — facial landmark regression.} King's Dlib library~\cite{king2009dlib}
popularised 68-point facial landmark detection via an ensemble of regression trees, enabling
reliable EAR computation~\cite{soukupova2016} without deep learning. Soukupov{\'a} and
{\v{C}}ech~\cite{soukupova2016} formalised EAR as a~single scalar for blink detection,
demonstrating frame-level detection accuracy above 99\,\% on controlled datasets. The
central limitation of this generation is the absence of a~WebAssembly-compatible
runtime, making client-side browser deployment impractical.

\emph{Third generation — deep landmark models on device.} MediaPipe
FaceLandmarker~\cite{lugaresi2019mediapipe} provides 468 normalised facial landmarks and
52 blendshape coefficients at approximately 30\,fps in a~browser WebAssembly context
using only a~CPU delegate. It subsumes and extends the Dlib approach: the denser
landmark set provides better coverage of the periorbital region, and the blendshape
scores offer a~calibrated signal for partial facial movements that EAR alone does not
capture. NeuroAssist AI builds on this generation, retaining EAR for the double-blink
state machine (where its simplicity aids interpretability) while using blendshape scores
for single-eye blink and eyebrow gestures.

\subsection{Deliberate Gesture Recognition and False-Positive Suppression}

A fundamental challenge in gesture-triggered AAC is distinguishing intentional
activations from the involuntary facial activity that continuously occurs during normal
behaviour. The spontaneous blink rate in healthy adults is 15--20 blinks per
minute~\cite{bentivoglio1997}; in MND patients, rates may be lower due to motor neuron
loss but remain non-zero. Any threshold-based approach applied at the single-frame level
therefore produces false-positive rates in the range of several per minute under natural
conditions, rendering unattended monitoring infeasible.

The two main strategies in the literature for suppressing false positives in facial AAC
are \emph{dwell time} and \emph{intentionality modelling}. Dwell-time selection — widely
used in gaze-based AAC — requires the user to fixate a~target region for a~fixed
duration (typically 0.5--1\,s)~\cite{majaranta2002eyetyping}. This approach is
inapplicable to blink-based input, as a~blink by definition closes the eye and disrupts
fixation. Intentionality modelling via state machines has been proposed in the broader
human-computer interaction literature~\cite{zander2011passive}, but its application to
blink-based AAC for MND has not been systematically described. The six-state machine
introduced in this paper (Section~\ref{sec:statemachine}) applies intentionality
constraints specific to the double-blink gesture: minimum closure duration per blink,
confirmed inter-blink opening, and a~temporal window — a~combination not found, to the
authors' knowledge, in prior AAC literature.

\subsection{Edge Computing and Privacy for Health Applications}

Processing health-related sensor data on-device rather than in the cloud has received
increasing attention following the enforcement of GDPR in 2018 and LGPD in
2020~\cite{edpb2020privacy}. Facial video is explicitly categorised as biometric data
under both frameworks, attracting the highest regulatory obligations for processing.
Local execution via WebAssembly provides a~practical privacy guarantee: the
specification's memory-isolated execution model prevents a~loaded module from accessing
any browser state outside its linear memory, and no network call is required to execute
the landmarker inference~\cite{haas2017wasm}. TensorFlow.js, ONNX Runtime Web, and
MediaPipe Tasks all leverage this runtime for on-device inference at speeds comparable
to native code on modern CPUs. Deployment latency of WebAssembly models for face
analysis has been reported at 20--40\,ms per frame on consumer hardware without GPU
acceleration~\cite{lugaresi2019mediapipe}, consistent with the 30\,fps operation
observed in our PoC.

\subsection{Software Architecture for Health Systems}

Hexagonal (Ports and Adapters) architecture~\cite{cockburn2005hexagonal} has been
proposed as a~structural pattern for health information systems where the domain model
must remain isolated from technology changes — a~recurring requirement in regulated
healthcare software where audit requirements, data schemas, and integration standards
evolve independently~\cite{richardson2018microservices}. The pattern's explicit separation
of inbound and outbound ports facilitates incremental extension without modifying core
clinical logic, a~property particularly valuable in assistive technology where new
gesture modalities are added as patients' motor capabilities change.

Microservice decomposition of AAC platforms has not been reported in the literature, to
the authors' knowledge; most existing systems are monolithic desktop or firmware
applications. The NeuroAssist AI architecture demonstrates that decomposing a
gesture-to-alert pipeline into independently deployable services (GEST, EVNT, and planned
PAC, ACS, HIST, DASH services) is viable at PoC scale, and that compile-time boundary
enforcement via ArchUnit provides a~practical mechanism for maintaining hexagonal
discipline in a~team-developed Java codebase.
